\slajdid{09_1-s029}
\begin{frame}[label=09_1-s029]
\gusttekst
Uslov zasićenja tranzistora u kolu je
\[V_{DSn,D}=V_O\geq\frac{(V_{GSn}-V_{Tn})}{1+\frac{(V_{GSn}-V_{Tn})}{L_nE_{Cn}}}=\frac{(V_I-V_{Tn})}{1+\frac{(V_I-V_{Tn})}{L_nE_{Cn}}}\]
Uočite da efekat zasićenja brzine nosilaca „brže“ uvodi tranzistor u zasićenje (pri manjim naponima $V_{DS}$) ali isto tako sa druge strane i „sporije“ izvodi tranzistor iz zasićenja.
\[v_O(t_{pHL})=\frac{V_{OH}+V_{OL}}2\geq\frac{V_{OHmin}-V_{Tn}}{1+\frac{(V_{OHmin}-V_{Tn})}{L_nE_{Cn}}}\]
\[\text{Tačno}\qquad\frac{V_{OH}+V_{OL}}2\geq\frac{L_nE_{Cn}(V_{OHmin}-V_{Tn})}{L_nE_{Cn}+(V_{OHmin}-V_{Tn})}\]
\[\text{Zanemarivanje}\quad\frac{V_{OH}+V_{OL}}2\geq L_nE_{Cn}\quad\text{ili}\quad\frac{V_{OH}+V_{OL}}2\geq(V_{OHmin}-V_{Tn})\]
Za standardne vrednosti biće zadovoljen uslov, tako da tranzistor radi u zasićenju za sve vreme ovog procesa pražnjenja kapacitivnosti na izlazu.
\end{frame}
